\documentclass[a4paper]{article}
\usepackage{graphicx}
\usepackage[top=10mm, bottom=20mm, left=30mm, right=30mm, includehead]{geometry}
\usepackage{fancyhdr}
\usepackage{hyperref}
\usepackage{mathptmx}
\usepackage{subcaption}
%\usepackage{parskip}
\usepackage{verbatim}
\usepackage{amsmath}
\usepackage{amssymb}
\usepackage{amsfonts}


\usepackage[backend=biber,style=authoryear]{biblatex}
\addbibresource{ref_DA.bib}

\pagestyle{fancy}
\fancyhf{}
\fancyfoot[C]{\thepage}
\renewcommand{\headrulewidth}{0pt}  % Elimina la línea en la cabecera
\date{}
\setlength{\headheight}{23.10004pt}
\setlength{\parindent}{0pt}


\fancypagestyle{plain}{
    \fancyhead[R]{\textit{Smart Sensing}}
    \fancyhead[L]{\textit{Ángel González Villatoro, Michael Hilfer, Paolo Olivucci}}
    \renewcommand{\headrulewidth}{0pt}
    }

\title{\bfseries\Large{Data Assimilation}}

\begin{document}
\maketitle 

\section{Introduction}

\subsection{Historical Development of Data Assimilation}

\textbf{1. Initial Motivation: Weather Prediction}\\

The concept of Data Assimilation emerged in the early 20th century with the attempt to use mathematical 
models for atmospheric prediction. Lewis Fry Richardson (1922) attempted to forecast weather by manually 
solving atmospheric equations \parencite{richardson1922weather}. However, the failure due to inconsistent 
initial conditions highlighted the need for a method to combine real-world data with models.\\

\textbf{2. Mathematical Formalization}\\

In the 1950s–60s, foundational algorithms were developed. Sasaki (1958) introduced the variational approach 
(3D-Var), minimizing the difference between model output and observations \parencite{sasaki1958objective}. 
Gandin (1963) then developed Optimal Interpolation (OI), a key method in operational meteorology 
\parencite{gandin1963objective}. These efforts introduced key concepts like background (prior forecast) and 
analysis (corrected state).\\

\textbf{3. Expansion with Supercomputing}\\

In the 1970s–80s, the rise of numerical weather models and observational data from satellites and radar 
made Data Assimilation essential. Institutions such as ECMWF and NCEP adopted it for global forecasting. 
The 4D-Var method also emerged, allowing time-evolving correction using forward and backward integrations.\\

\textbf{4. Stochastic Methods}\\

In the 1990s–2000s, ensemble-based methods like the Ensemble Kalman Filter (EnKF) by Evensen and Houtekamer 
became popular. These are more robust for nonlinear systems and allow estimation of uncertainty by running 
multiple perturbed simulations.\\

\textbf{5. Cross-Disciplinary Expansion}\\

Today, DA is applied in diverse fields:
\begin{itemize}
    \item \textbf{Oceanography and climatology}
    \item \textbf{Aerodynamics and fluid mechanics}
    \item \textbf{Cardiovascular modeling} \parencite{bertoglio2012identification}
    \item \textbf{Wildfire forecasting} \parencite{rochoux2020towards}
    \item \textbf{Energy systems and engineering control} \parencite{saredi2022state}
    \item \textbf{Medicine and oncology} \parencite{maier2021bayesian}
    \item \textbf{Economics and uncertainty quantification} \parencite{wikle2007bayesian}
\end{itemize}

\subsection{Relationships to Other Key Concepts}

\textbf{Design of Experiments (DoE)}: DA benefits from well-designed measurement strategies and can guide 
them in return. DoE optimizes when and where to measure to extract maximum information. DA enables adaptive 
observation targeting (e.g., in OSSEs), improving future experiment planning \parencite{kalnay2003atmospheric}.\\

\textbf{Predictive Uncertainty}: DA is not only about estimating system states, but also about their 
associated uncertainty. Methods like EnKF and Bayesian DA provide probabilistic forecasts, crucial in 
fields like meteorology and medicine \parencite{evensen1994sequential, wikle2007bayesian}.\\

\textbf{Mode Decomposition}: Techniques like POD, DMD, and SVD are often used in tandem with DA. Modal 
decomposition reduces model order, making DA computationally feasible and efficient. Alternatively, DA 
can reconstruct modal coefficients from partial data \parencite{benner2015survey}.\\

\textbf{Machine Learning (ML)}: DA is increasingly integrated with ML. Hybrid models leverage DA for 
physical structure and ML for flexibility. Physics-Informed Neural Networks (PINNs) incorporate model physics 
into neural training, while DA can refine internal states or assist in training under noise \parencite{raissi2019physics, benner2015survey}.

\section{Key Components of Data Assimilation}

Data Assimilation (DA) integrates model predictions with observations to estimate the evolving state of a physical 
system. Regardless of the specific method (e.g., variational, ensemble-based, hybrid), most DA frameworks share the 
following fundamental components \parencite{kalnay2003atmospheric, evensen1994sequential, wikle2007bayesian}.\\

\textbf{1. Model (Forward Operator)}\\

The model is typically a set of dynamical equations—often derived from physics—that describe the time evolution of 
the system's state. This could be:\\

\begin{itemize}
  \item Navier–Stokes equations for fluid flow,
  \item Atmospheric or oceanic circulation models,
  \item Biomechanical or epidemiological models.
\end{itemize}

The model provides the \textit{background forecast} (also known as the \textit{prior}) used in the assimilation 
process \parencite{kalnay2003atmospheric}.\\

\textbf{2. Observations (Data)}\\

These are real-world measurements, often noisy and partial, obtained from:
\begin{itemize}
  \item Sensors (e.g., weather stations, PIV cameras),
  \item Satellites, drones, or in-situ probes,
  \item Medical imaging or industrial diagnostics.
\end{itemize}

DA seeks to correct the model using these observations while respecting the system dynamics \parencite{wikle2007bayesian}.\\

\textbf{3. Observation Operator (H)}\\

This maps the model state into the observational space \parencite{law2015data}. It may be:\\

\begin{itemize}
  \item A simple projection (e.g., extracting velocity at a point),
  \item A complex nonlinear operator (e.g., computing a Doppler signal or pressure field).
\end{itemize}

This operator allows comparing model predictions to actual measurements.\\

\textbf{4. Error Statistics (Covariances)}\\

DA explicitly or implicitly uses error models for:\\

\begin{itemize}
  \item Background error (uncertainty in the model prediction),
  \item Observation error (noise, bias, or measurement limitations).
\end{itemize}

These are often encoded in covariance matrices:\\

\begin{itemize}
  \item $\mathbf{B}$ for background,
  \item $\mathbf{R}$ for observations.
\end{itemize}

In ensemble methods, these are estimated from sample statistics \parencite{evensen1994sequential}.\\

\textbf{5. Assimilation Algorithm}\\

This is the mathematical method that combines model and data. Examples include:\\

\begin{itemize}
  \item Variational methods: 3D-Var, 4D-Var,
  \item Sequential methods: Kalman Filter, EnKF,
  \item Bayesian or particle-based methods.
\end{itemize}

The goal is to find an \textit{analysis state}, which is the best estimate of the true system, balancing fidelity 
to the model and the observations \parencite{kalnay2003atmospheric}.\\

\textbf{6. Update Cycle (Sequential or Smoothing)}\\

In \textit{filtering}, assimilation is performed each time new data arrives (e.g., at every timestep). In 
\textit{smoothing}, future data is used to update past states. In both cases, assimilation is often repeated 
iteratively over a time window \parencite{law2015data}.\\

\begin{table}[h!]
\centering
\caption{Summary of Key Components}
\begin{tabular}{|l|p{10cm}|}
\hline
\textbf{Component} & \textbf{Description} \\
\hline
Model (M) & Governs system evolution \\
Observations (y) & Real-world data inputs \\
Observation Operator (H) & Projects model to data space \\
Covariance Matrices (B, R) & Represent uncertainties in model and measurements \\
Assimilation Algorithm & Merges data and model (e.g., EnKF, 4D-Var) \\
Update Mechanism & Filter (online) or smoother (retrospective) \\
\hline
\end{tabular}
\end{table}

\section{Methods of Data Assimilation}

\subsection{Nudging}

Nudging, also known as \textit{Newtonian relaxation}, is one of the earliest and simplest data assimilation 
techniques. It works by adding a correction term to the model dynamics, proportional to the difference between 
observations and the model state \parencite{stauffer1993use, kalnay2003atmospheric}:\\

\[
\frac{dx}{dt} = f(x) + \alpha \left( y - H(x) \right)
\]

Here:\\

\begin{itemize}
  \item $x$ is the system state,
  \item $f(x)$ is the model evolution,
  \item $y$ is the observation,
  \item $H(x)$ maps the model to observation space,
  \item $\alpha$ is the nudging (relaxation) coefficient.
\end{itemize}

\textbf{Advantages:}\\

\begin{itemize}
  \item Simple to implement and computationally efficient.
  \item Useful for stabilizing or guiding model trajectories.
\end{itemize}

\textbf{Limitations:}\\

\begin{itemize}
  \item No uncertainty quantification.
  \item Sensitive to the choice of $\alpha$.
  \item Ignores spatial correlations in errors.
\end{itemize}

\begin{comment}
\textbf{Example — Simple: Heat diffusion in a room}\\

Consider a 1D heat equation:
\[
\frac{\partial T}{\partial t} = \kappa \frac{\partial^2 T}{\partial x^2}
\]

Given a single temperature observation at $x_0$, we apply nudging:
\[
\frac{\partial T}{\partial t} = \kappa \frac{\partial^2 T}{\partial x^2} + \alpha \left( T_{obs} - T(x_0, t) \right) \delta(x - x_0)
\]

This forces the temperature at $x_0$ to approach the observed value while still evolving according to the diffusion 
model.\\


\textbf{Example — Complex: Global atmospheric models}\\

Before the adoption of 4D-Var and EnKF, nudging was widely used by operational centers like ECMWF and GFDL to 
assimilate temperature and wind data into global models. For instance, the ERA-40 reanalysis used nudging to 
maintain large-scale consistency with observations \parencite{uppala2005era}.\\

A more advanced form is \textit{spectral nudging}, where only specific spatial scales (e.g., synoptic structures) 
are nudged, allowing finer scales to evolve freely.\\
\end{comment}

\begin{comment}
\subsubsection{Simple Nudging Example: 1D Heat Diffusion with a Single Observation}

\textbf{Physical Model}\\

Consider the temperature evolution in a 1D rod of length $L$ governed by the heat equation:
\[
\frac{\partial T(x,t)}{\partial t} = \kappa \frac{\partial^2 T(x,t)}{\partial x^2}
\]
with homogeneous boundary conditions:
\[
T(0,t) = T(L,t) = 0, \quad T(x,0) = T_0(x)
\]

\textbf{Observation Setup}\\

Assume we have one temperature sensor located at position $x = x_m$ that records the true temperature $T_{\text{obs}}(t) \approx T(x_m,t)$ over time.\\

\textbf{Nudging-Augmented Model}\\

To incorporate the observation, we modify the equation by adding a nudging term:
\[
\frac{\partial T(x,t)}{\partial t} = \kappa \frac{\partial^2 T(x,t)}{\partial x^2} + \alpha (T_{\text{obs}}(t) - T(x,t)) \delta(x - x_m)
\]
where:
\begin{itemize}
  \item $\alpha$ is the nudging gain,
  \item $\delta(x - x_m)$ is a Dirac delta function centered at the sensor location.
\end{itemize}

\subsection*{Purpose and Interpretation}

This example demonstrates:
\begin{itemize}
  \item How localized observations influence a dynamical system,
  \item The impact of the nudging gain $\alpha$,
  \item How nudging helps reconcile simulated and observed data,
  \item That even a single observation can correct the full state over time via diffusion.
\end{itemize}

This is a pedagogical test case ideal for teaching nudging principles.



\subsubsection{Application Example: Turbulent Flow Around a Square Cylinder}

Zauner et al. (2020) apply nudging to a URANS simulation of flow around a square cylinder at $Re = 22{,}000$. DNS data is used as observation input to guide the URANS solution.

\subsubsection*{Model Setup}

\begin{itemize}
  \item Governing equations: URANS with $k$--$\omega$ SST model.
  \item Geometry: 2D square cylinder in cross-flow.
  \item Observations: velocity fields from DNS.
  \item Nudging applied: in a spatial region downstream of the cylinder, with $\alpha = 1$.
\end{itemize}

\subsubsection*{Equation}

Nudged momentum equation:
\[
\frac{\partial \mathbf{u}}{\partial t} + (\mathbf{u} \cdot \nabla) \mathbf{u} = -\nabla p + \nu \nabla^2 \mathbf{u} + \alpha (\mathbf{u}_{\text{DNS}} - \mathbf{u})
\]

\subsubsection*{Results}

\begin{itemize}
  \item Vortex shedding frequency ($St \approx 0.137$) was accurately recovered.
  \item High-frequency KH instabilities ($St \approx 4.384$) were also captured.
  \item SPOD analysis showed spectral agreement between nudged URANS and DNS.
\end{itemize}

\textbf{Conclusion:} Nudging allowed URANS to capture multiscale flow features with minimal complexity, bridging the gap between low-cost and high-fidelity simulations.

\end{comment}

\textbf{Modern variants:}\\

\begin{itemize}
  \item \textit{Adaptive nudging:} The $\alpha$ coefficient varies with time or observation quality.
  \item \textit{Physics-informed learning:} Nudging principles are sometimes embedded in machine learning models to enforce physical consistency.
\end{itemize}


\subsection{Kalman Filter (KF)}

The Kalman Filter (KF) is a recursive data assimilation algorithm that provides the optimal estimate of the state 
of a linear dynamical system with Gaussian errors. It updates both the state estimate and its associated uncertainty 
as new observations become available \parencite{kalman1960new, kalnay2003atmospheric}.\\

\textbf{System model:}\\

\begin{align}
x_{k} &= A x_{k-1} + w_{k-1}, \quad w_{k-1} \sim \mathcal{N}(0, Q) \\
y_{k} &= H x_{k} + v_{k}, \quad v_{k} \sim \mathcal{N}(0, R)
\end{align}

Where:\\

\begin{itemize}
  \item $x_k$ is the system state at time $k$,
  \item $A$ is the system transition matrix,
  \item $H$ is the observation operator,
  \item $Q$ and $R$ are the process and observation error covariances.
\end{itemize}

\textbf{KF Equations:}\\

\begin{itemize}
  \item \textit{Prediction step:}
    \[
    \hat{x}_{k}^{-} = A \hat{x}_{k-1}, \quad P_k^{-} = A P_{k-1} A^T + Q
    \]
  \item \textit{Update step:}
    \[
    K_k = P_k^{-} H^T (H P_k^{-} H^T + R)^{-1}
    \]
    \[
    \hat{x}_{k} = \hat{x}_k^{-} + K_k (y_k - H \hat{x}_k^{-}), \quad P_k = (I - K_k H) P_k^{-}
    \]
\end{itemize}

\textbf{Advantages:}\\

\begin{itemize}
  \item Provides statistically optimal estimate under Gaussian linear assumptions.
  \item Updates uncertainty along with state.
  \item Recursive and efficient for low-dimensional systems.
\end{itemize}

\textbf{Limitations:}\\

\begin{itemize}
  \item Requires linear dynamics and observation models.
  \item Assumes Gaussian errors.
  \item Poor performance in nonlinear or high-dimensional systems.
\end{itemize}

\begin{comment}

\textbf{Example — Simple: GPS tracking of a vehicle}\\

Consider tracking the position and velocity of a car in 1D using GPS (position measurements) and an accelerometer 
(model-based prediction). The Kalman Filter fuses noisy measurements with a constant-acceleration motion model to 
provide smoother estimates and quantify uncertainty.\\

\textbf{Example — Complex: Real-time atmospheric estimation}\\

In meteorology, early efforts used the Kalman Filter for small regional weather models to assimilate pressure and 
wind observations. However, due to computational costs, it was replaced by EnKF for high-dimensional systems. Still, 
it is used in subsystems such as satellite orbit tracking or in embedded filters in robotics and navigation systems.\\

\textbf{Modern developments:}\\

\begin{itemize}
  \item \textit{Extended Kalman Filter (EKF)}: handles mild nonlinearity using linearization.
  \item \textit{Unscented Kalman Filter (UKF)}: uses deterministic sampling for better performance in nonlinear systems.
  \item \textit{Ensemble Kalman Filter (EnKF)}: approximates KF using Monte Carlo ensemble statistics (covered in next section).
\end{itemize}

\end{comment}

\subsection{Ensemble Kalman Filter (EnKF)}

The Ensemble Kalman Filter (EnKF) extends the classical Kalman Filter to nonlinear and high-dimensional systems 
using a Monte Carlo approach. Instead of computing and propagating the full covariance matrix, it uses an ensemble 
of state realizations to approximate forecast and analysis error statistics \parencite{evensen1994sequential, houtekamer2005ensemble}.\\

\textbf{System model:}\\

\begin{align}
x_k^{(i)} &= \mathcal{M}(x_{k-1}^{(i)}) + w_{k-1}^{(i)}, \quad i = 1, ..., N_e \\
y_k &= H x_k + v_k
\end{align}

Where:\\

\begin{itemize}
  \item $x_k^{(i)}$: the $i$-th ensemble member at time $k$,
  \item $\mathcal{M}$: nonlinear model operator,
  \item $H$: (possibly linear or nonlinear) observation operator,
  \item $w_k^{(i)}$, $v_k$: random process and observation noise.
\end{itemize}

\textbf{Basic algorithm steps:}\\

\begin{enumerate}
  \item \textit{Forecast step:} Propagate each ensemble member using the model.
  \item \textit{Analysis step:} Compute sample mean and covariances:
    \[
    \bar{x}_f = \frac{1}{N_e} \sum_{i=1}^{N_e} x_k^{(i)}, \quad
    P_f = \frac{1}{N_e - 1} \sum_{i=1}^{N_e} (x_k^{(i)} - \bar{x}_f)(x_k^{(i)} - \bar{x}_f)^T
    \]
  \item Update each member using:
    \[
    x_k^{(i,a)} = x_k^{(i,f)} + K \left( y_k + \eta_k^{(i)} - H x_k^{(i,f)} \right)
    \]
    where $K$ is the Kalman gain computed from the ensemble statistics and $\eta_k^{(i)}$ is perturbed observation noise.
\end{enumerate}

\textbf{Advantages:}\\

\begin{itemize}
  \item Suitable for nonlinear and high-dimensional systems.
  \item No need to compute or store full covariance matrices.
  \item Provides flow-dependent error covariances.
\end{itemize}

\textbf{Limitations:}\\

\begin{itemize}
  \item Sampling errors for small ensembles.
  \item Requires ensemble inflation and localization to control spurious correlations.
\end{itemize}

\begin{comment}
\textbf{Example — Simple: Hydrologic water level estimation}\\

A rainfall-runoff model (e.g., for a reservoir) is simulated with 20 ensemble members, each with different initial 
soil moisture. Streamflow observations are assimilated to correct the ensemble and improve flood forecasting.\\

\textbf{Example — Complex: Wildfire spread modeling}\\

\textcite{rochoux2020towards} apply EnKF with a surrogate model based on Polynomial Chaos to predict wildfire 
spread. Observations from thermal imagery and wind conditions are used to adjust uncertain model parameters in 
real time.\\

\subsubsection{Kalman Filter (EnKF): Hurricane and Wildfire Modeling}

\textbf{Hurricane Katrina (Torn and Hakim)} \parencite{torn2009ensemble}:  
EnKF was applied to real observations from the 2005 Hurricane Katrina event. The study assimilated data from dropsondes, satellite radiances, and surface stations into the WRF model. Key outcomes:
\begin{itemize}
  \item Improved track and intensity forecasts over deterministic methods.
  \item Demonstrated benefits of ensemble-derived flow-dependent covariances.
  \item Showed that assimilation frequency and observation type strongly influence forecast quality.
\end{itemize}

\textbf{Wildfire Spread (Rochoux et al.)} \parencite{rochoux2020towards}:  
A reduced-order surrogate model based on Polynomial Chaos Expansion (PCE) was coupled with EnKF to estimate key parameters of wildfire spread in real time. Main features:
\begin{itemize}
  \item Thermal imagery and meteorological data used as observations.
  \item Demonstrated orders-of-magnitude faster runtime compared to full physics-based models.
  \item Effective estimation of uncertain inputs like fuel moisture and wind speed.
\end{itemize}

\end{comment}

\textbf{Modern applications:}\\

\begin{itemize}
  \item Operational oceanographic and meteorological forecasting (e.g., NCEP).
  \item Reservoir engineering (e.g., oil/gas field assimilation).
  \item Environmental hazard prediction (e.g., wildfires, pollution spread).
\end{itemize}


\subsection{Three-Dimensional Variational Data Assimilation (3D-Var)}

The 3D-Var method is a widely used variational approach in data assimilation. It estimates the state of a system 
by minimizing a cost function that balances the discrepancy between the model background (prior) and the 
observations \parencite{lorenc1986analysis, courtier1994strategy}.\\

\textbf{Cost function:}\\

\[
J(x) = \frac{1}{2}(x - x_b)^T B^{-1} (x - x_b) + \frac{1}{2}(y - H(x))^T R^{-1} (y - H(x))
\]

Where:\\

\begin{itemize}
  \item $x$: analysis state (to be estimated),
  \item $x_b$: background state (model forecast),
  \item $y$: observations,
  \item $H(x)$: observation operator (possibly nonlinear),
  \item $B$: background error covariance matrix,
  \item $R$: observation error covariance matrix.
\end{itemize}

The solution $x_a$ minimizes $J(x)$ and represents the best estimate combining model and observation information.\\

\textbf{Advantages:}\\

\begin{itemize}
  \item Robust and computationally efficient.
  \item Suitable for large-scale operational systems.
  \item Well-established theoretical and numerical framework.
\end{itemize}

\textbf{Limitations:}\\

\begin{itemize}
  \item Only uses observations at a single time (no time window).
  \item Requires careful specification of $B$ and $R$.
  \item Assumes stationary background statistics.
\end{itemize}

\begin{comment}

\textbf{Example — Simple: Flow estimation in a wind tunnel}\\

A steady 2D incompressible flow is simulated over a flat plate. Sparse pressure measurements from the surface are 
assimilated using 3D-Var to estimate the full velocity field. The background is provided by a coarse CFD simulation, 
and $H$ maps states to surface pressures.\\

\textbf{Example — Complex: Operational weather forecasting}\\

Centers like Météo-France and NCEP used 3D-Var operationally for years. For instance, in NCEP's Global Data 
Assimilation System (GDAS), surface and satellite observations are assimilated every 6 hours to produce an analysis 
state used for forecasting. It replaced Optimal Interpolation and later evolved into hybrid 3DEnVar systems.\\

\textbf{Implementation considerations:}\\

\begin{itemize}
  \item Minimization is typically performed via iterative gradient-based methods.
  \item Preconditioning and multigrid approaches help accelerate convergence.
  \item Hybrid versions combine ensemble-based flow-dependent covariances with static $B$ matrices.
\end{itemize}

\end{comment}

\subsection{Four-Dimensional Variational Data Assimilation (4D-Var)}

The 4D-Var method extends 3D-Var by incorporating time-evolving observations over a specified assimilation window. 
It minimizes a cost function over the initial condition $x_0$ such that the model trajectory fits both the 
background and all observations within the time window \parencite{courtier1994strategy, rabier2000ecmwf}.\\

\textbf{Cost function:}\\

\[
J(x_0) = \frac{1}{2}(x_0 - x_b)^T B^{-1} (x_0 - x_b) + \frac{1}{2} \sum_{i=0}^{N} (y_i - H_i(x_i))^T R^{-1} (y_i - H_i(x_i))
\]

Where:\\

\begin{itemize}
  \item $x_0$: initial state to optimize,
  \item $x_b$: background state (prior),
  \item $y_i$: observations at time $t_i$,
  \item $x_i = \mathcal{M}_{0 \rightarrow i}(x_0)$: model forecast from $x_0$ to $t_i$,
  \item $H_i$: observation operator at $t_i$,
  \item $B$: background covariance, $R$: observation covariance.
\end{itemize}

\textbf{Key features:}\\

\begin{itemize}
  \item Uses time-distributed observations, increasing observational impact.
  \item Requires tangent linear and adjoint models for gradient computation.
\end{itemize}

\textbf{Advantages:}\\

\begin{itemize}
  \item Accurate use of dynamics and time-correlated observations.
  \item Strong theoretical basis in control theory and inverse problems.
  \item Proven performance in operational weather forecasting.
\end{itemize}

\textbf{Limitations:}\\

\begin{itemize}
  \item Requires development of adjoint and linearized models.
  \item Computationally expensive, especially for large windows.
  \item Assumes linearity of model error within the window.
\end{itemize}

\begin{comment}

\textbf{Example — Simple: 1D heat equation inversion}\\

Given temperature observations at several times and locations in a rod, 4D-Var can estimate the initial temperature 
distribution that leads to best fit over time using forward integration of the heat equation and backward adjoint 
optimization.\\

\textbf{Example — Complex: ECMWF global atmospheric analysis}\\

The ECMWF uses 4D-Var to assimilate millions of observations from satellites, aircraft, balloons, and surface 
stations into a high-resolution global atmospheric model. The assimilation window spans 6–12 hours and includes 
nonlinear dynamics, satellite radiative transfer, and complex observation operators \parencite{rabier2000ecmwf}.\\

\textbf{Variants and developments:}\\

\begin{itemize}
  \item \textit{Incremental 4D-Var:} solves in reduced space using linearized operators.
  \item \textit{Hybrid 4DEnVar:} combines 4D-Var structure with ensemble-based covariances.
  \item \textit{Weak-constraint 4D-Var:} accounts for model errors explicitly.
\end{itemize}

\subsection{Variational Methods (4D-Var): RANS Mean Flow Reconstruction}

Foures et al. \parencite{foures2014data} developed a 4D-Var technique to reconstruct Reynolds-averaged Navier–Stokes (RANS) mean flows. DNS data of flow around a circular cylinder were subsampled to simulate sparse observations. Highlights include:
\begin{itemize}
  \item Reconstruction of velocity fields and Reynolds stress distributions from partial data.
  \item Use of a steady adjoint solver to minimize a cost function defined over time-averaged states.
  \item Demonstrated the ability to extrapolate information into unmeasured regions.
\end{itemize}

Extra (D. Rival)\\

\textbf{State Observer / PID-Inspired Assimilation Application: Pressure-Driven Flow (Neetson et al.)}\\

Neetson and Rival introduced a state observer method inspired by PID control, where 
assimilation is applied in the pressure Poisson equation instead of momentum equations. 
The method dynamically corrected both pressure and velocity fields and was robust even 
under noisy conditions. The assimilation was tested in flow with vortex shedding, showing 
improved frequency response and stability.

\end{comment}

\section{Mathematical Implementation of Data Assimilation Methods}

This section presents the detailed mathematical formulation and step-by-step computational implementation of the most widely used data assimilation (DA) techniques.

\subsection{Kalman Filter (KF)}

\textbf{Goal:} Estimate the state $\hat{x}_k$ of a linear dynamical system at time $k$ using a recursive filter.

\textbf{Model:}
\[
x_k = A x_{k-1} + w_{k-1}, \quad w_{k-1} \sim \mathcal{N}(0, Q)
\]
\[
y_k = H x_k + v_k, \quad v_k \sim \mathcal{N}(0, R)
\]

\textbf{Procedure:}
\begin{itemize}
  \item \textbf{Initialization:} Provide initial guess $\hat{x}_0$, $P_0$.
  \item \textbf{Forecast:}
  \[
  \hat{x}_k^{-} = A \hat{x}_{k-1}, \quad P_k^{-} = A P_{k-1} A^T + Q
  \]
  \item \textbf{Update:}
  \[
  K_k = P_k^{-} H^T (H P_k^{-} H^T + R)^{-1}
  \]
  \[
  \hat{x}_k = \hat{x}_k^{-} + K_k (y_k - H \hat{x}_k^{-}), \quad P_k = (I - K_k H) P_k^{-}
  \]
\end{itemize}

\subsection{Ensemble Kalman Filter (EnKF)}

\textbf{Goal:} Extend Kalman filter for nonlinear models using ensemble statistics.

\textbf{Forecast:} For each ensemble member $i$:
\[
x_k^{(i,f)} = \mathcal{M}(x_{k-1}^{(i)})
\]

\textbf{Statistics:}
\[
\bar{x}_f = \frac{1}{N_e} \sum x_k^{(i,f)}, \quad
P_f = \frac{1}{N_e - 1} \sum (x_k^{(i,f)} - \bar{x}_f)(x_k^{(i,f)} - \bar{x}_f)^T
\]

\textbf{Update:}
\[
K_k = P_f H^T (H P_f H^T + R)^{-1}
\]
\[
x_k^{(i,a)} = x_k^{(i,f)} + K_k \left(y_k + \eta_k^{(i)} - H x_k^{(i,f)} \right)
\]

\subsection{Three-Dimensional Variational Method (3D-Var)}

\textbf{Goal:} Find the best state $x$ that minimizes a cost function at a single time.

\[
J(x) = \frac{1}{2}(x - x_b)^T B^{-1}(x - x_b) + \frac{1}{2}(y - H(x))^T R^{-1} (y - H(x))
\]

\textbf{Solution:} Use optimization (e.g., gradient descent) to find $x_a = \arg\min J(x)$.

\subsection{Four-Dimensional Variational Method (4D-Var)}

\textbf{Goal:} Optimize initial condition $x_0$ to minimize mismatch between model trajectory and observations over a time window.

\[
J(x_0) = \frac{1}{2}(x_0 - x_b)^T B^{-1}(x_0 - x_b) + \frac{1}{2} \sum_{i=0}^{N} (y_i - H_i(x_i))^T R^{-1} (y_i - H_i(x_i))
\]

\textbf{Where:} $x_i = \mathcal{M}_{0 \rightarrow i}(x_0)$

\textbf{Procedure:}
\begin{itemize}
  \item Use forward model to propagate $x_0$ and calculate residuals.
  \item Use adjoint model to compute gradient of $J$.
  \item Apply iterative minimization to update $x_0$.
\end{itemize}

\subsection{Nudging}

\textbf{Goal:} Introduce a feedback correction term into the model equations.

\[
\frac{dx}{dt} = f(x) + \alpha (y - H(x))
\]

\textbf{Implementation:}
\begin{itemize}
  \item At each time step, compute the discrepancy $e = y - H(x)$.
  \item Add a relaxation term $\alpha e$ to guide the state toward the observation.
  \item $\alpha$ controls the strength of the correction.
\end{itemize}


\section{Nudging: Theory and Application}

\subsection{Overview of the Nudging Method}

Nudging, also known as Newtonian relaxation, is a data assimilation (DA) technique that introduces a feedback term into the governing equations of a dynamical system to guide its evolution toward observational data. It is particularly effective in scenarios where real-time data assimilation is needed with low computational cost.

The general idea is to augment the model equations with a relaxation term that penalizes the difference between the model state and available measurements.

\subsubsection*{Mathematical Formulation}

Consider a generic dynamical system:
\[
\frac{dx}{dt} = f(x)
\]

In nudging, this becomes:
\[
\frac{dx}{dt} = f(x) + \alpha (y - H(x))
\]

Where:
\begin{itemize}
  \item $x$: model state vector,
  \item $f(x)$: model dynamics,
  \item $y$: vector of observations,
  \item $H(x)$: observation operator projecting model state to measurement space,
  \item $\alpha$: nudging coefficient controlling the strength of feedback.
\end{itemize}

The correction term $ \alpha (y - H(x)) $ adjusts the model state toward the observations over time.

\subsubsection*{Interpretation}

- If $\alpha$ is large, the model rapidly follows observations (but risks instability).
- If $\alpha$ is too small, the influence of observations is negligible.
- $H(x)$ can be linear (e.g., extracting a value at a location) or nonlinear.

Nudging can also be formulated in continuous or discrete time and extended to include spatial operators or localized feedback.

\subsection{Application to Turbulent Flow Around a Square Cylinder}

Zauner et al. \parencite{zauner2020nudging} applied nudging to the turbulent flow past a square cylinder at a Reynolds number $Re = 22{,}000$ using a URANS solver.

\subsubsection*{Governing Equations}

The nudged momentum equation in their study is written as:
\[
\frac{\partial \bar{u}_\alpha}{\partial t} + (\bar{u}_\alpha \cdot \nabla)\bar{u}_\alpha + \nabla \bar{p}_\alpha - \nabla \cdot \left[2\left(\frac{1}{Re} + \nu_t(\tilde{\nu}_\alpha)\right) \nabla_s \bar{u}_\alpha\right] = \alpha H^\dagger [H(\bar{u}_\alpha) - m(t)]
\]

Where:
\begin{itemize}
  \item $\bar{u}_\alpha$: nudged velocity field,
  \item $\bar{p}_\alpha$: pressure,
  \item $\nu_t$: eddy viscosity from Spalart–Allmaras model,
  \item $H$: operator extracting observed velocity values,
  \item $H^\dagger$: adjoint operator redistributing the mismatch to the model space,
  \item $m(t)$: observed (reference) velocity data from DNS.
\end{itemize}

\subsubsection*{Implementation Steps}

\begin{enumerate}
  \item Synthetic reference data were generated from DNS of the same flow.
  \item Observations were spatially averaged in the spanwise direction and provided at uniform grid points.
  \item Nudging was applied in a subdomain downstream of the cylinder, at every timestep.
  \item The nudging coefficient $\alpha$ was tuned for stability and effectiveness.
\end{enumerate}

\subsubsection*{Key Results}

\begin{itemize}
  \item \textbf{Low-frequency (vortex shedding)}: The Strouhal number of vortex shedding was corrected from $St = 0.126$ (URANS) to $St = 0.137$ (DNS-like) using nudging.
  \item \textbf{High-frequency (Kelvin–Helmholtz)}: Successfully recovered small-scale instabilities with sufficient spatial and temporal resolution in the nudged data.
  \item \textbf{Mean flow accuracy}: Time-averaged quantities such as recirculation zone and pressure distribution improved significantly.
\end{itemize}

\subsubsection*{Interpretation}

The nudging method effectively synchronizes URANS simulations with high-fidelity DNS data, allowing reconstruction of flow structures across different scales. The feedback mechanism corrects both phase and amplitude of the dynamic structures, even outside the nudged region.



\section{Kalman Filter: Theory and Applications}

\subsection{Theoretical Foundations}

The Kalman Filter (KF), introduced by Rudolf E. Kalman in 1960 \parencite{kalman1960new}, is a recursive 
algorithm that provides the optimal estimate of a system’s internal state based on a series of incomplete and 
noisy measurements. It is especially powerful for linear systems with Gaussian noise.

\subsubsection*{State-Space Model}

A discrete-time linear system is modeled as:
\[
x_k = A x_{k-1} + B u_{k-1} + w_{k-1}
\]
\[
z_k = H x_k + v_k
\]

Where:
\begin{itemize}
  \item $x_k$: system state at time $k$,
  \item $u_k$: control input,
  \item $z_k$: observation at time $k$,
  \item $A$: state transition matrix,
  \item $B$: control matrix,
  \item $H$: observation matrix,
  \item $w_k$, $v_k$: zero-mean Gaussian process and observation noise, with covariances $Q$ and $R$.
\end{itemize}

\subsubsection*{Kalman Filter Algorithm}

\textbf{Prediction step:}
\[
\hat{x}_{k|k-1} = A \hat{x}_{k-1|k-1} + B u_{k-1}
\]
\[
P_{k|k-1} = A P_{k-1|k-1} A^T + Q
\]

\textbf{Update step:}
\[
K_k = P_{k|k-1} H^T (H P_{k|k-1} H^T + R)^{-1}
\]
\[
\hat{x}_{k|k} = \hat{x}_{k|k-1} + K_k (z_k - H \hat{x}_{k|k-1})
\]
\[
P_{k|k} = (I - K_k H) P_{k|k-1}
\]

\subsection{Applications from Literature}

\subsubsection*{Example 1: Target Tracking System}

In \textcite{ghanai2009kalman}, the KF is adapted as a \textit{Kalman Controller} to guide a mobile platform (e.g., laser or antenna) to track a moving target in azimuth and elevation. The radar provides periodic updates on the target's location. The KF estimates control actions to realign the platform, even in the presence of noise and disturbances.

Mathematically, the controller is based on a 1D Kalman Filter with the state representing control actions and the measurements as positional errors. Control is applied via DC motors for azimuth and rise.

Key features:
\begin{itemize}
  \item Recursive estimation avoids storage of entire trajectory.
  \item Resilient to white noise and load disturbances.
  \item Demonstrated strong tracking and disturbance rejection.
\end{itemize}

\subsubsection*{Example 2: Fuzzy Model Identification}

In another application, the KF is employed to estimate parameters in a Takagi–Sugeno fuzzy model for nonlinear system identification, such as ECG signal approximation \parencite{ghanai2009kalman}. The model has two estimation stages:

\begin{enumerate}
  \item \textbf{Premise function estimation:} Using fuzzy clustering and Kalman filtering to identify membership functions.
  \item \textbf{Consequent parameter estimation:} Using a separate KF (KF2) for linear regression over fuzzy rule outputs.
\end{enumerate}

This method showed strong performance for noisy and nonlinear time series.


\section{Ensemble Kalman Filter: Theory and Applications}

\subsection{Theoretical Overview}

The Ensemble Kalman Filter (EnKF) is an extension of the Kalman Filter designed to handle nonlinear and high-dimensional systems. Originally proposed by Geir Evensen \parencite{evensen1994sequential}, it replaces the explicit computation of error covariance matrices with Monte Carlo estimates using an ensemble of model realizations.

\subsubsection*{State Estimation Framework}

Assume a dynamical model:
\[
x_k^{(i)} = \mathcal{M}(x_{k-1}^{(i)}) + w_{k-1}^{(i)}, \quad i = 1, \ldots, N
\]

And observations:
\[
y_k = H x_k + v_k
\]

Where:
\begin{itemize}
  \item $x_k^{(i)}$: $i$-th ensemble member at time $k$,
  \item $\mathcal{M}$: nonlinear model operator,
  \item $H$: observation operator (linear or nonlinear),
  \item $w_k$, $v_k$: process and observation noise.
\end{itemize}

\subsubsection*{Forecast Step}

Each ensemble member is propagated forward in time:
\[
x_k^{(i,f)} = \mathcal{M}(x_{k-1}^{(i)})
\]

Compute sample mean and covariance:
\[
\bar{x}_f = \frac{1}{N} \sum x_k^{(i,f)}, \quad
P_f = \frac{1}{N - 1} \sum (x_k^{(i,f)} - \bar{x}_f)(x_k^{(i,f)} - \bar{x}_f)^T
\]

\subsubsection*{Analysis Step}

Perturb observations (optional) and compute the Kalman gain:
\[
K_k = P_f H^T (H P_f H^T + R)^{-1}
\]

Update each ensemble member:
\[
x_k^{(i,a)} = x_k^{(i,f)} + K_k \left(y_k + \eta_k^{(i)} - H x_k^{(i,f)}\right)
\]
Where $\eta_k^{(i)} \sim \mathcal{N}(0, R)$.

\subsection{Applications from Literature}

\subsubsection*{Example 1: Quasi-Geostrophic Ocean Model (Evensen, 1994)}

Evensen applied EnKF to a nonlinear barotropic quasi-geostrophic model of ocean circulation. The filter assimilated sparse observations and successfully reconstructed evolving vorticity fields. Key contributions:
\begin{itemize}
  \item Demonstrated robustness for non-Gaussian distributions.
  \item Introduced covariance localization and inflation strategies.
  \item Laid the groundwork for sequential DA in geophysical models.
\end{itemize}

\subsubsection*{Example 2: Atmospheric EnKF for Operational Forecasting (Houtekamer et al.)}

In \textcite{houtekamer2005ensemble}, the EnKF was implemented in the Canadian global forecasting system. Innovations included:
\begin{itemize}
  \item Localization to reduce spurious long-distance correlations.
  \item Adaptive inflation to stabilize ensemble spread.
  \item Assimilation of millions of observations including satellite radiances.
\end{itemize}


\subsubsection*{Example 3: Wildfire Spread Modeling (Rochoux et al.)}

\textcite{rochoux2020towards} employed an Ensemble Kalman Filter framework enhanced with a Polynomial Chaos Expansion (PCE) surrogate model to reduce computational cost in wildfire spread forecasting. Their methodology includes:

\begin{itemize}
  \item Construction of a surrogate model using PCE to approximate the expensive forward wildfire simulation (FIREFLY).
  \item Use of EnKF to assimilate thermal camera observations to estimate uncertain parameters (e.g., wind, vegetation moisture).
  \item Implementation of a reduced-cost DA loop capable of providing timely forecasts.
\end{itemize}

This demonstrated that EnKF is feasible for large-scale operational use and provided comparable results to 4D-Var.

\section{Ensemble Kalman Filter: Theory and Application to Wildfire Forecasting}

\subsection{Theoretical Overview}

The Ensemble Kalman Filter (EnKF) is a Monte Carlo-based sequential data assimilation technique that uses an ensemble of model realizations to estimate the state and error statistics of nonlinear, high-dimensional systems \parencite{evensen1994sequential}. The ensemble approach avoids the need to evolve full covariance matrices explicitly.

\subsubsection*{Formulation}

Consider a nonlinear model:
\[
x_k^{(i)} = \mathcal{M}(x_{k-1}^{(i)}) + w_{k-1}^{(i)}
\]
\[
y_k = \mathcal{H}(x_k) + v_k
\]

where:
\begin{itemize}
  \item $x_k^{(i)}$: state vector of the $i$-th ensemble member,
  \item $\mathcal{M}$: nonlinear model operator,
  \item $\mathcal{H}$: observation operator,
  \item $w_k^{(i)}$, $v_k$: process and observation noise.
\end{itemize}

The forecast ensemble yields mean $\bar{x}_f$ and covariance $P_f$. Each member is updated as:
\[
x_k^{(i,a)} = x_k^{(i,f)} + K_k \left( y_k + \eta_k^{(i)} - \mathcal{H}(x_k^{(i,f)}) \right)
\]
with Kalman gain:
\[
K_k = P^{f}_{k} H^T (H P^{f}_{k} H^T + R)^{-1}
\]

\subsection{Application: Wildfire Spread Forecasting (Rochoux et al.)}

In \textcite{rochoux2020towards}, the EnKF is applied to the simulation of wildfire spread using a reduced-cost surrogate model based on Polynomial Chaos Expansion (PCE). The study integrates thermal camera observations to infer uncertain environmental parameters affecting fire behavior.

\subsubsection*{Model and Parameters}

The FIREFLY model is used to simulate the wildfire front. Uncertain parameters include:
\begin{itemize}
  \item Fuel moisture content,
  \item Surface-to-volume ratio of grass fuel,
  \item Wind speed and direction.
\end{itemize}

\subsubsection*{Observation Operator}

The observation operator projects the model state (fire front) into a space where distances from observed points can be computed:
\[
d_t = [(x_1 - x_1^o), (x_2 - x_2^o), \ldots] \in \mathbb{R}^{2N_o}
\]

\subsubsection*{EnKF Setup}

\begin{itemize}
  \item An ensemble of model runs is generated using samples of input parameters.
  \item A PCE surrogate approximates the output of the forward model, significantly reducing cost.
  \item The EnKF is used to update the estimated parameters by assimilating fire front positions derived from thermal imagery.
\end{itemize}

\subsubsection*{Outcomes}

\begin{itemize}
  \item Real-time correction of model forecasts using sparse data.
  \item Improved tracking of the evolving fire front compared to open-loop predictions.
  \item A 10-fold reduction in computational cost compared to the original model.
\end{itemize}



\section{Variational Methods: 3D-Var and 4D-Var}

Variational data assimilation methods estimate the state of a dynamical system by minimizing a cost function that measures the discrepancy between the model state and available observations, weighted by their error covariances. They rely on optimization rather than sequential updates.

\subsection{3D-Var: Three-Dimensional Variational Assimilation}

\subsubsection*{Formulation}

3D-Var seeks the analysis state $x_a$ that minimizes:
\[
J(x) = \frac{1}{2}(x - x_b)^T B^{-1} (x - x_b) + \frac{1}{2}(y - H(x))^T R^{-1} (y - H(x))
\]

Where:
\begin{itemize}
  \item $x_b$: background (prior) state,
  \item $y$: observation vector,
  \item $H$: observation operator,
  \item $B$: background error covariance,
  \item $R$: observation error covariance.
\end{itemize}

\textbf{Solution:} Use iterative methods (e.g., conjugate gradients) to minimize $J(x)$.

\subsubsection*{Application: Weather Prediction Systems}

As discussed in \textcite{navon2009data}, 3D-Var was operationally implemented at major centers (ECMWF, Météo-France) for atmospheric analysis. Observations from various sources (surface, upper-air, satellites) are assimilated every 6 hours to initialize forecasts.

\textbf{Advantages:}
\begin{itemize}
  \item Stable and computationally efficient.
  \item Easy to implement with static background covariances.
\end{itemize}

\textbf{Limitations:}
\begin{itemize}
  \item Observations are used at a single time only.
  \item Assumes static error statistics.
\end{itemize}

\subsection{4D-Var: Four-Dimensional Variational Assimilation}

\subsubsection*{Formulation}

4D-Var extends 3D-Var by assimilating observations distributed over a time window $[t_0, t_N]$. It minimizes:
\[
J(x_0) = \frac{1}{2}(x_0 - x_b)^T B^{-1}(x_0 - x_b) + \frac{1}{2} \sum_{i=0}^{N} \left( y_i - H_i(x_i) \right)^T R^{-1} \left( y_i - H_i(x_i) \right)
\]

with:
\[
x_i = \mathcal{M}_{0 \rightarrow i}(x_0)
\]

Where $\mathcal{M}$ is the nonlinear model mapping initial state to time $t_i$.

\textbf{Gradient:}
The minimization requires the gradient of $J$ with respect to $x_0$, which is computed via the adjoint of the tangent linear model:
\[
\nabla J = B^{-1}(x_0 - x_b) + \sum_{i=0}^{N} \mathcal{M}^*_{i \rightarrow 0} H_i^T R^{-1}(H_i(x_i) - y_i)
\]

\subsubsection*{Applications}

In \textcite{courtier1994strategy} and \textcite{courtier1999assimilation}, 4D-Var was implemented for global weather forecasting at ECMWF. It assimilated satellite radiances, aircraft, and balloon observations to improve initial condition accuracy over 12-hour windows.

\textbf{Advantages:}
\begin{itemize}
  \item Dynamical consistency: integrates model dynamics directly.
  \item Efficient use of time-distributed data.
\end{itemize}

\textbf{Limitations:}
\begin{itemize}
  \item Requires adjoint and tangent linear models.
  \item High computational cost.
\end{itemize}


\section{Detailed Mathematical Description of Variational Data Assimilation Methods}

Variational data assimilation (Var-DA) techniques seek to find the optimal estimate of the system state by minimizing a cost function that combines information from a background forecast and observations. These methods are rooted in control theory and inverse problems, and they form the theoretical basis for high-performance data assimilation systems.

\subsection{3D-Var (Three-Dimensional Variational Assimilation)}

\subsubsection*{Mathematical Formulation}

The goal of 3D-Var is to find the analysis state $x_a$ that minimizes the following cost function:

\[
J(x) = \frac{1}{2}(x - x_b)^T B^{-1}(x - x_b) + \frac{1}{2}(y - H(x))^T R^{-1}(y - H(x))
\]

Where:
\begin{itemize}
  \item $x$: model state vector (control variable),
  \item $x_b$: background state (typically a short-range forecast),
  \item $y$: observation vector,
  \item $H$: observation operator mapping model space to observation space,
  \item $B$: background error covariance matrix,
  \item $R$: observation error covariance matrix.
\end{itemize}

\textbf{Minimization Process:}
\begin{enumerate}
  \item Define $J(x)$ and compute its gradient $\nabla J$:
  \[
  \nabla J = B^{-1}(x - x_b) - H^T R^{-1}(y - H(x))
  \]
  \item Use an iterative optimization algorithm (e.g., quasi-Newton, conjugate gradient).
  \item Solve in the control variable space or reduced control space via preconditioning.
\end{enumerate}

\subsubsection*{Application: Operational Weather Forecasting (Navon, Bouttier \& Courtier)}

3D-Var was a cornerstone in early operational weather centers like ECMWF and Météo-France \parencite{navon2009data, courtier1999assimilation}. It assimilated diverse observations such as radiosondes, surface stations, and satellite radiances every 6 hours to produce an initial state for numerical weather prediction models.

\subsection{4D-Var (Four-Dimensional Variational Assimilation)}

\subsubsection*{Mathematical Formulation}

4D-Var extends 3D-Var by incorporating observations over a time window $[t_0, t_N]$, optimizing the initial condition $x_0$ such that the model trajectory best fits time-distributed observations:

\[
J(x_0) = \frac{1}{2}(x_0 - x_b)^T B^{-1}(x_0 - x_b) + \frac{1}{2} \sum_{i=0}^{N} \left( y_i - H_i(x_i) \right)^T R^{-1} \left( y_i - H_i(x_i) \right)
\]

Subject to:
\[
x_i = \mathcal{M}_{0 \rightarrow i}(x_0)
\]

Where $\mathcal{M}$ is the nonlinear model operator advancing the state forward in time.

\textbf{Gradient Computation:}
The gradient of $J$ is computed using the adjoint of the tangent linear model:
\[
\nabla J = B^{-1}(x_0 - x_b) + \sum_{i=0}^{N} \mathcal{M}^*_{i \rightarrow 0} H_i^T R^{-1}(H_i(x_i) - y_i)
\]

This requires:
\begin{itemize}
  \item A tangent linear model (TLM) $\mathcal{M}'$,
  \item An adjoint model $\mathcal{M}^*$ for back-propagation of sensitivity information.
\end{itemize}

\textbf{Optimization:} As with 3D-Var, iterative gradient-based algorithms (e.g., L-BFGS) are used to find $x_0^*$.

\subsubsection*{Application: ECMWF 4D-Var System (Courtier et al.)}

In \textcite{courtier1994strategy}, the authors developed and implemented the incremental 4D-Var algorithm at ECMWF. It used a sequence of linearized inner loops within a nonlinear outer loop, improving both convergence and computational feasibility.

\textbf{Implementation Features:}
\begin{itemize}
  \item Assimilation window of 12 hours.
  \item Use of TL/adjoint versions of the IFS (Integrated Forecasting System) model.
  \item Inner loops use a simplified version of the dynamics to accelerate convergence.
\end{itemize}

\subsubsection*{Advantages and Limitations}

\textbf{Advantages:}
\begin{itemize}
  \item Time consistency: observations at multiple times are fully utilized.
  \item Physically consistent trajectory constrained by model dynamics.
  \item Accurate use of correlated observations (e.g., satellite radiances).
\end{itemize}

\textbf{Limitations:}
\begin{itemize}
  \item High computational cost due to model integrations and adjoint calls.
  \item Complexity in generating tangent linear and adjoint models.
\end{itemize}


% === Extended Variational Methods Section ===
\section{Detailed Theory of Variational Methods}

Variational data assimilation methods aim to determine the most probable state of a physical system by minimizing a cost function that incorporates prior information (background) and observational data. These techniques leverage optimization frameworks and require a solid understanding of error covariance modeling and adjoint-based gradient computation.

\subsection{3D-Var: Three-Dimensional Variational Assimilation}

3D-Var minimizes a cost function of the form:
\[
J(x) = \frac{1}{2}(x - x_b)^T B^{-1}(x - x_b) + \frac{1}{2}(y - H(x))^T R^{-1}(y - H(x))
\]
where $x$ is the model state, $x_b$ the background (forecast), $y$ the observations, $B$ the background error covariance matrix, $R$ the observation error covariance matrix, and $H$ the observation operator (possibly nonlinear).

The minimum of $J(x)$ is found using iterative solvers, such as conjugate gradient methods. Preconditioning and variable transformations may be applied to improve convergence, especially in high-dimensional systems.

\textbf{Example: Météo-France operational 3D-Var system} assimilates surface pressure, temperature, and satellite radiances. The model state includes wind, temperature, and humidity. $B$ is constructed from climatological statistics and balance operators.

\subsection{4D-Var: Four-Dimensional Variational Assimilation}

4D-Var extends the 3D-Var framework by incorporating observations across a time window $[t_0, t_N]$. The objective is to optimize the initial state $x_0$ so that the model trajectory $\mathcal{M}(x_0)$ fits observations at multiple times:

\[
J(x_0) = \frac{1}{2}(x_0 - x_b)^T B^{-1}(x_0 - x_b) + \frac{1}{2} \sum_{i=0}^{N} (y_i - H_i(x_i))^T R^{-1} (y_i - H_i(x_i))
\]
with $x_i = \mathcal{M}_{0 \rightarrow i}(x_0)$ where $\mathcal{M}$ is the numerical forecast model.

To compute the gradient $\nabla J$, adjoint and tangent linear models are used:
\[
\nabla J = B^{-1}(x_0 - x_b) + \sum_{i=0}^N \mathcal{M}^*_{i \rightarrow 0} H_i^T R^{-1} (H_i(x_i) - y_i)
\]

\textbf{Example: ECMWF’s Incremental 4D-Var} (Courtier et al., 1994) implemented an efficient scheme using outer loops (nonlinear model) and inner loops (linearized system) for faster convergence. Assimilation windows were 12 hours, with millions of observations assimilated.

\textbf{Example: Foures et al. (2014)} applied 4D-Var to reconstruct RANS-averaged flow fields. The cost function was based on mean velocity from DNS, and adjoint models of the RANS equations were derived. The process allowed estimation of unobserved flow components constrained by physics.



% === Extended Descriptions of Sequential Methods ===

\section{Detailed Theory of Sequential Methods}

Sequential DA methods estimate the state recursively as new data become available. These include filtering techniques that predict forward in time and update when observations arrive.

\subsection{Nudging}

Nudging, or Newtonian relaxation, modifies model equations to draw the state towards observations:
\[
\frac{dx}{dt} = f(x) + \alpha(y - H(x))
\]
Here, $\alpha$ is a gain matrix or scalar that determines the influence of observations.

\textbf{Zauner et al. (2020):} Nudging was applied to a turbulent flow simulation around a square cylinder using URANS. DNS data served as synthetic observations. The nudging term was added only in the near wake, improving vortex shedding and turbulence structure. Sensitivity to temporal and spatial sampling was analyzed.

\subsection{Kalman Filter}

For linear systems with Gaussian noise:
\[
x_k = A x_{k-1} + w_k, \quad y_k = H x_k + v_k
\]
KF updates the state mean and covariance via:
\begin{align*}
\hat{x}_k^- &= A \hat{x}_{k-1}, \\
P_k^- &= A P_{k-1} A^T + Q, \\
K_k &= P_k^- H^T (H P_k^- H^T + R)^{-1}, \\
\hat{x}_k &= \hat{x}_k^- + K_k (y_k - H \hat{x}_k^-), \\
P_k &= (I - K_k H) P_k^-
\end{align*}

\textbf{Ghanai and Chafaa:} KF was used in two applications:
\begin{itemize}
  \item Control of radar antennas: estimated position and velocity of targets under white noise and disturbance.
  \item Nonlinear ECG modeling: Used two-stage KF to update both fuzzy rule base (premise) and consequent parameters in a Takagi–Sugeno model.
\end{itemize}

\subsection{Ensemble Kalman Filter (EnKF)}

EnKF is a Monte Carlo approximation of KF:
\begin{align*}
x_k^{(i,f)} &= \mathcal{M}(x_{k-1}^{(i)}) + w_k^{(i)} \\
K_k &= P^f H^T (H P^f H^T + R)^{-1} \\
x_k^{(i,a)} &= x_k^{(i,f)} + K_k(y_k + \eta_k^{(i)} - H x_k^{(i,f)})
\end{align*}
Where $P^f$ is the sample covariance of the ensemble.

\textbf{Evensen (1994):} Applied EnKF to ocean modeling. Used 50-member ensemble in a QG model with sparse observations to track vorticity fields.

\textbf{Torn and Hakim (2009):} Applied EnKF with WRF for real-time forecasting of Hurricane Katrina. Assimilated aircraft and radar observations. Showed sensitivity to ensemble spread and boundary condition perturbations.

\textbf{Rochoux et al. (2020):} Applied EnKF to wildfire modeling. Combined it with a Polynomial Chaos surrogate to reduce model cost. Fire front positions were assimilated from thermal imagery, updating uncertain parameters such as wind and fuel moisture.



% === Detailed Theory and Applications of Data Assimilation Methods ===

\section{Mathematical Formulation of Data Assimilation}

Let $x_k$ be the state vector, $x_b$ the background (prior), $y_k$ the observations, $H$ the observation operator, $B$ the background error covariance matrix, and $R$ the observation error covariance matrix. The general objective is to obtain the best estimate $x_a$ by combining the prior model forecast and available data, accounting for their respective uncertainties.

\section{Sequential Data Assimilation Methods}

These methods update the state estimate recursively over time as new observations become available.

\subsection{Nudging}

Nudging adds a relaxation term to the model equations:
\[
\frac{dx}{dt} = f(x) + \alpha(y - H(x))
\]
The coefficient $\alpha$ controls the strength of the relaxation toward the observed state.

\textbf{Application: Zauner et al. (2020)}\\
In a study of turbulent flow around a square cylinder, nudging was applied in a URANS framework. Synthetic DNS observations were used to correct large-scale vortex shedding. Sensitivity analysis revealed optimal nudging performance depends on the spatial and temporal resolution of observations.

\subsection{Kalman Filter (KF)}

For linear-Gaussian systems:
\[
x_k = A x_{k-1} + w_k, \quad y_k = H x_k + v_k
\]
With recursive steps:
\begin{align*}
\text{Predict:} & \quad \hat{x}_k^- = A \hat{x}_{k-1}, \quad P_k^- = A P_{k-1} A^T + Q \\
\text{Update:} & \quad K_k = P_k^- H^T (H P_k^- H^T + R)^{-1} \\
& \quad \hat{x}_k = \hat{x}_k^- + K_k (y_k - H \hat{x}_k^-), \quad P_k = (I - K_k H) P_k^-
\end{align*}

\textbf{Applications:}
\begin{itemize}
  \item \textbf{Ghanai and Chafaa:} KF used for radar antenna tracking and fuzzy ECG modeling. Dual Kalman Filters estimated both system states and parameters of nonlinear models.
\end{itemize}

\subsection{Ensemble Kalman Filter (EnKF)}

EnKF uses an ensemble of state realizations:
\begin{align*}
x_k^{(i,f)} &= \mathcal{M}(x_{k-1}^{(i)}) + w_k^{(i)} \\
K_k &= P^f H^T (H P^f H^T + R)^{-1} \\
x_k^{(i,a)} &= x_k^{(i,f)} + K_k(y_k + \eta_k^{(i)} - H x_k^{(i,f)})
\end{align*}
Where $P^f$ is the forecast ensemble covariance.

\textbf{Applications:}
\begin{itemize}
  \item \textbf{Evensen (1994):} Demonstrated EnKF in oceanography using a QG model.
  \item \textbf{Torn and Hakim (2009):} Applied EnKF to Hurricane Katrina with the WRF model. Assessed ensemble perturbation strategies.
  \item \textbf{Rochoux et al. (2020):} Used EnKF with a Polynomial Chaos surrogate model for wildfire spread prediction. Fire front locations were corrected using thermal camera observations.
\end{itemize}

\section{Variational Data Assimilation Methods}

Variational methods find the optimal state by minimizing a cost function using all available observations at once (3D-Var) or over a time window (4D-Var).

\subsection{3D-Var: Three-Dimensional Variational Assimilation}

\[
J(x) = \frac{1}{2}(x - x_b)^T B^{-1}(x - x_b) + \frac{1}{2}(y - H(x))^T R^{-1}(y - H(x))
\]

Gradient descent is used to solve:
\[
\nabla J = B^{-1}(x - x_b) - H^T R^{-1}(y - H(x))
\]

\textbf{Application: Bouttier and Courtier (1999)}\\
3D-Var was implemented at Météo-France to assimilate surface, upper-air, and satellite data. $B$ was derived from climatological statistics and balance constraints.

\subsection{4D-Var: Four-Dimensional Variational Assimilation}

4D-Var extends 3D-Var by assimilating data over time:
\[
J(x_0) = \frac{1}{2}(x_0 - x_b)^T B^{-1}(x_0 - x_b) + \frac{1}{2} \sum_{i=0}^{N} (y_i - H_i(x_i))^T R^{-1} (y_i - H_i(x_i))
\]
Where $x_i = \mathcal{M}_{0 \rightarrow i}(x_0)$. Gradients use adjoint models:
\[
\nabla J = B^{-1}(x_0 - x_b) + \sum_{i=0}^{N} \mathcal{M}^*_{i \rightarrow 0} H_i^T R^{-1}(H_i(x_i) - y_i)
\]

\textbf{Applications:}
\begin{itemize}
  \item \textbf{Courtier et al. (1994):} ECMWF’s 4D-Var system used nested loops with linearized dynamics for operational weather prediction.
  \item \textbf{Foures et al. (2014):} Applied 4D-Var to reconstruct mean RANS flow fields using sparse PIV data and adjoint RANS solvers.
\end{itemize}

\printbibliography

\end{document}
